\chapter{权 1 形式的 Galois 表示}
% ===== PDF p.079 =====
\noindent\textbf{（Ben Heuer 撰）}

\noindent\textbf{上次}：对权 $\ge 2$ 的新形式，我们构造了附着其上的 Galois 表示：

\noindent\textbf{定理 10.1}（Deligne，1971）\phantomsection\label{thm:10.1}\quad 设 $N\ge 3$、$k\ge 2$，$\varepsilon:(\Z/N\Z)^\times\to\C^\times$ 为 Dirichlet 特征。设 $f\in M_k(\Gamma_0(N),\varepsilon)$ 是权 $k$ 的新形式，$K$ 为由诸 $a_p(f)$ 与 $\varepsilon$ 生成的数域。则对任一素数 $\ell\nmid N$ 以及 $K$ 中位于 $\ell$ 之上的位 $\lambda$，存在连续、奇、不可约的 Galois 表示
\[
\rho:G_\Q\longrightarrow\GL_2(K_\lambda),
\]
它在 $N$ 外非分歧，且对一切 $p\nmid N\ell$，$\Frob_p$ 的特征多项式为
\[
X^2-a_pX+\varepsilon(p)p^{k-1}.
\]

\subsection*{附着于权 1 模形式的 Galois 表示}

\begin{itemize}
\item 今日目标：证明下面的定理。
\end{itemize}

\noindent\textbf{定理 10.2}（Deligne--Serre，1973）\phantomsection\label{thm:10.2}\quad 设 $N\ge 3$，$\varepsilon:(\Z/N\Z)^\times\to\C^\times$ 为 Dirichlet 特征，$f\in M_1(\Gamma_0(N),\varepsilon)$ 是权 $1$ 的新形式。则存在连续、奇、不可约的 Galois 表示
\[
\rho:G_\Q\longrightarrow\GL_2(\C),
\]
它在 $N$ 外非分歧，且对一切 $p\nmid N$，$\Frob_p$ 的特征多项式为
\[
X^2-a_pX+\varepsilon(p).
\]

\begin{itemize}
\item 至此完成了「新形式 $\rightsquigarrow$ Galois 表示」的构造。
\item 注意 $\rho$ 的像落在 $\GL_2(\C)$ 中，即 $\rho$ 是 Artin 表示；这与权 $\ge 2$ 的情形相反，那里得到的是 $\ell$-进表示。为什么？
\begin{itemize}
\item 简单事实：任何 Artin 表示的像都有限。
\item 事实上，$\rho$ 已经可以在数域上定义，因此也可以像前面那样嵌入 $p$-进域 $K_\lambda$；但用 $\C$ 是为了标示「像有限」。
\item 从 Galois 侧看，$\rho$ 只有在 $k=1$ 时才可能像有限：因为要求 $\det\rho=\varepsilon\chi_{\mathrm{cycl}}^{k-1}$，而 $\chi_{\mathrm{cycl}}$ 的像无限。
\end{itemize}
\item 若 $\varepsilon$ 为偶，则 $M_1(\Gamma_0(N),\varepsilon)=0$。因此 $\varepsilon$ 为奇，从而 $\rho$ 为奇。
\end{itemize}

\subsection*{构造的第一版草图}

与上次相比，这个证明的几何成分更少，Galois 表示论的成分更多：

\begin{itemize}
\item 把 $f$ 模 $p$ 约化，得到一个「模 $p$ 模形式」。
\item 利用模形式的同余，说明存在模 $p$ 的 Hecke 本征形式 $f'$，它与 $f$ 有相同的特征值，但权更高，$\ge 2$。
\item 用「Deligne--Serre 提升\idx{Deligne--Serre 提升}」把 $f'$ 提升回特征 $0$。
\item 用权 $\ge 2$ 的构造，附着 Galois 表示 $\rho_{f'}$。
\item 把 Galois 表示模 $\ell$ 约化。
\item 把 Galois 表示从 $\F_\ell$ 提升到 $\C$。
\end{itemize}
% ===== PDF p.080 =====
\subsection*{模 $p$ 模形式初识}

\begin{itemize}
\item 设 $N\ge 3$，$S$ 为任一 $\Z[1/N]$-概形。
\item 回忆：我们曾把 $S$ 上权 $k$ 的模形式定义为
\[
M_k(\Gamma_1(N);S):=\Gamma\bigl(X^{*}_{\Gamma_1(N),S},\omega^k\bigr).
\]
\item $\omega$ 是万有椭圆曲线上相对微分的丰沛线丛。
\item 可以定义 $q$-展开：设 $S=\Spec(A)$，则在 $\Z((q))\otimes_\Z A$ 上有 Tate 曲线 $T(q)$；在 $\Spec(\Z((q))\otimes_\Z A)$ 上，有典范同构
\[
\omega_{\Gm}=\omega_{T(q)}.
\]
\item 线丛 $\omega_{\Gm}$ 平凡：$\Gm$ 上的典范微分 $\mathrm{d}x/x$ 是可逆截面。
\item 借助这个平凡化，对 $X^{*}_{\Gamma_1(N),A}$ 的每个尖点都得到相应的 $q$-展开
\[
M_k(\Gamma_1(N);S)\longrightarrow A[[q]].
\]
\end{itemize}

\noindent\textbf{定理 10.3}（Katz 的 $q$-展开原理）\phantomsection\label{thm:10.3}\quad 对每个尖点、所有 $k$，该映射都是单射。

\begin{itemize}
\item 现在把这一切应用于 $S=\Spec(F)$，其中 $F$ 是特征 $p$ 的有限域：所得的模形式称为\textbf{模 $p$ 模形式}。
\item 下面的事实或许有助于体会模 $p$ 模形式：
\end{itemize}

\noindent\textbf{引理 10.4}\phantomsection\label{lem:10.4}\quad 对权 $k\ge 2$，任何模 $p$ 模形式都是某个数域上的模形式模 $p$ 的约化。特别地，此时
\[
M_k(\Gamma_1(N);F)=M_k\bigl(\Gamma_1(N);\Z[1/N]\bigr)\otimes_\Z F.
\]

\noindent\textbf{证明}\quad 用 Riemann--Roch 可以证明 $H^1\bigl(X^{*}_{\Gamma_1(N)},\omega^{\otimes k}\bigr)=0$。\hfill$\square$

\begin{itemize}
\item 然而，通常给出更内禀的定义会更好。例如：
\end{itemize}

\subsection*{模 $p$ 模形式的例子：Hasse 不变量\idx{Hasse 不变量}}

\begin{itemize}
\item 回忆：对特征 $p$ 的任一椭圆曲线，有 Verschiebung 同源 $V:E^{(p)}\to E$。
\item 特别地，对万有椭圆曲线也有。
\item 在相对微分上，这给出反方向的拉回映射 $V^*:\omega\to\omega^{\otimes p}$。
\item 与 $\omega^{-1}$ 张量，得
\[
V^*:\OO\longrightarrow\omega^{\otimes(p-1)},\qquad\text{即}\quad
V^*\in\Gamma\bigl(\omega^{\otimes(p-1)}\bigr)=M_{p-1}(\Gamma_1(N);F).
\]
\item 这就是 \textbf{Hasse 不变量}，记作 $\mathrm{Ha}$。
\item 它是权 $p-1$ 的模 $p$ 模形式。事实上：它是级 $1$ 的本征形式。
\end{itemize}

\noindent\textbf{命题 10.5}\phantomsection\label{prop:10.5}\quad Hasse 不变量的 $q$-展开为常数 $1\in\F_p[[q]]$。

\noindent\textbf{证明（略述）}\quad 利用 $\omega_{T(q)}=\omega_{\Gm}$：在 $\Gm$ 上有 $F=[p]$，从而 $V=1$。\hfill$\square$

\noindent\textbf{注}：$1\in M_0(\Gamma_1(N);\F_p)$ 也是 $q$-展开为 $1$ 的模 $p$ 模形式。

\noindent\textbf{注}：这与 $q$-展开原理并不矛盾，因为两者的权不同。\znote{标准定义如下：对特征 $p$ 的椭圆曲线 $E$ 与非零微分 $\eta$，由 $F^*\eta=A(E,\eta)\eta$ 得到权 $p-1$ 的模形式 $A$，其 $q$-展开恒为 $1$；$A$ 的零点恰为 $E$ 的超奇异点（即 $E[p]$ 为 $\alpha_p$ 的扩张之处）。因此乘以 $\mathrm{Ha}$ 的映射在超奇异点处退化，这正是要挖去这些点、转向 $p$-进模形式理论的原因。[来源：B 级讲义 \url{https://math.mit.edu/nt/Murilo2020-02-26.pdf}（2026-11-26 访问）。]}\par
% ===== PDF p.081 =====
\noindent\textbf{要点}：在特征 $p$ 下出现了一个新现象：

\begin{itemize}
\item 不同权的模形式可以有相同的 $q$-展开！
\item 例如，对任一权 $k$ 的模形式 $f$ 与任一 $n\in\N$，$\mathrm{Ha}^n\cdot f$ 是权 $k+n(p-1)$ 的模形式，且与 $f$ 有相同的 $q$-展开。
\item 特别地，这保持本征形式：Hecke 算子的同一特征值系统可以对应不同权的本征形式。
\item 乘法映射
\[
M_k(\Gamma_1(N);S)\xrightarrow{\ \cdot\,\mathrm{Ha}\ }M_{k+p-1}(\Gamma_1(N);S)
\]
是单射（$q$-展开原理），但不是满射！
\item 原因：$\mathrm{Ha}$ 有零点，且其零点恰为 $X^{*}_{\Gamma_1(N)}(\Fq)$ 的超奇异点。
\item 这是挖去这些点的动机之一；挖去之后便得到 $p$-进模形式的理论。
\end{itemize}

\subsection*{Eisenstein 级数}

\begin{itemize}
\item 对此还有一个更经典的视角，我们至少应当提一下：
\end{itemize}

\noindent\textbf{定义 10.6}\phantomsection\label{def:10.6}\quad 对任一偶 $k\ge 4$，令
\[
G_k:=\sum_{\substack{(n,m)\in\Z^2\\(n,m)\neq(0,0)}}\frac{1}{(m+n\tau)^k}.
\]
这是权 $k$、级 $1$ 的模本征形式，其 $q$-展开为
\[
2\zeta(k)\Bigl(1-\frac{2k}{B_k}\sum_{n=1}^{\infty}\sigma_{k-1}(n)q^n\Bigr),
\]
其中 $\sigma_{k-1}(n)=\sum_{d\mid n}d^{k-1}$，$B_k$ 是第 $k$ 个 Bernoulli 数。把归一化的 Eisenstein 级数定义为
\[
E_k:=\frac{1}{2\zeta(k)}G_k=1-\frac{2k}{B_k}\sum_{n=1}^{\infty}\sigma_{k-1}(n)q^n\in\Z_{(p)}[[q]].
\]
\znoteplain{原文（提取文本）此处作 $\frac{4}{B_k}$；按 $G_k$ 的标准 Fourier 展开应为 $\frac{2k}{B_k}$。核对：$k=4$ 时 $-\frac{2k}{B_4}=-\frac{8}{-1/30}=240$，正是 $E_4=1+240q+\cdots$ 的系数。译文按 $\frac{2k}{B_k}$ 处理。}

\noindent\textbf{引理 10.7}\phantomsection\label{lem:10.7}\quad 设 $p\ge 5$。则就 $q$-展开而言，有
\[
E_{p-1}\equiv 1\pmod p.
\]

\noindent\textbf{证明}\quad 这由 Bernoulli 数的 Kummer 同余得到。\hfill$\square$

\noindent\textbf{要点}：当 $p\ge 5$ 时，Eisenstein 级数 $E_{p-1}$ 模 $p$ 的约化就是 $\mathrm{Ha}$：
\[
E_{p-1}\equiv\mathrm{Ha}\pmod p,
\]
此即「Deligne 同余」。

\noindent\textbf{策略}（Deligne--Serre）：

\begin{itemize}
\item 给定权 $1$ 的本征形式，模 $p$ 约化。
\item 乘以 $\mathrm{Ha}$，使其成为权 $\ge 2$ 的本征形式。
\item 在此情形使用上一讲的构造。
\end{itemize}

\subsection*{附着于模 $p$ 模形式的 Galois 表示}
% ===== PDF p.082 =====
\noindent\textbf{定理 10.8}\phantomsection\label{thm:10.8}\quad 设 $k\ge 1$，$f\in M_k(\Gamma_1(N),\varepsilon;\F_p)$ 是模 $p$ 的尖点本征形式。设 $F_f$ 为由诸 $a_\ell(f)$ 与 $\varepsilon$ 在 $\F_p$ 上生成的（有限）域。则存在半单 Galois 表示
\[
\rho_f:G_\Q\longrightarrow\GL_2(F_f),
\]
它在 $Np$ 外非分歧，且对一切 $\ell\nmid Np$，
\[
\operatorname{charpoly}(\Frob_\ell\mid\rho_f)=X^2-a_\ell X+\varepsilon(\ell)\ell^{k-1}.
\]
\znote{原文此处 $\ell$ 与 $p$ 混用：先说「在 $N\ell$ 外非分歧」，又说「对一切 $\ell\nmid Np$」，且把特征多项式写成 $\Frob_p$ 的。按上下文，$p$ 是系数域 $\F_p$ 的特征，变动的素数应记作 $\ell$；译文统一为「在 $Np$ 外非分歧，且对一切 $\ell\nmid Np$」。相应地，$F_f$ 由诸 $a_\ell(f)$ 与 $\varepsilon$ 生成。}

\noindent\textbf{第一步：归约到 $k\ge 2$ 的情形。}若 $k=1$，乘以 $\mathrm{Ha}$。这不改变 $q$-展开，因此仍是本征形式，但其权 $\ge 2$。

\noindent\textbf{第二步：把 Hecke 特征系统提升到特征 $0$。}当权 $\ge 2$ 时这总是可行的，下面就来讨论。

\subsection*{Deligne--Serre 提升引理}

\begin{itemize}
\item 设 $R$ 是无零因子的 $\Z[1/N]$-代数。设
\[
\mathbb{T}_R\subseteq\End\bigl(M_k(\Gamma_1(N);R)\bigr)
\]
为 \textbf{Hecke 代数}，即由 Hecke 算子 $T_l,S_l$ 生成的 $R$-子代数。
\item \textbf{Hecke 特征系统}是环同态 $\psi:\mathbb{T}_R\to R$。Hecke 特征系统的典型例子如下：
\item 设 $g\in M_k(\Gamma_1(N);R)$ 是本征形式，它给出 Hecke 特征系统
\[
\psi_g:\mathbb{T}_R\longrightarrow R,\qquad T\longmapsto a_T,\quad\text{使得}\ T(g)=a_Tg.
\]
\end{itemize}

\noindent\textbf{引理 10.9}\phantomsection\label{lem:10.9}\quad 设 $R=K$ 是域。则任一 Hecke 特征系统 $\psi:\mathbb{T}_K\to K$ 都来自唯一一个 Hecke 本征形式。

\noindent\textbf{证明}。
\begin{itemize}
\item $\mathbb{T}_K$ 是 Artinian $K$-代数。
\item 取扩域后，它是 Artinian 局部环 $\mathbb{T}_i$ 的乘积。
\item 由于 $\mathbb{T}_K$ 忠实作用在 $M:=M_k(\Gamma_1(N);K)$ 上，子模 $M_i:=\mathbb{T}_i\cdot M$ 非零（交换代数事实）。
\item 在 $M_i$ 中任取本征向量，它的特征系统就是 $\psi$。
\item 唯一性：Hecke 特征值决定 $q$-展开。\hfill$\square$
\end{itemize}

\begin{itemize}
\item 设 $K$ 为数域，$\p$ 为 $K$ 中位于 $p\nmid N$ 之上的位，$\F_{\p}$ 为剩余域。
\item 记 $\OO_{K,(\p)}$ 为 $K$ 中 $\p$-整元素构成的赋值子环。
\end{itemize}

\noindent\textbf{引理 10.10}（Deligne--Serre 提升引理）\phantomsection\label{lem:10.10}\quad 设 $k\ge 2$，$g\in M_k(\Gamma_1(N);\F_{\p})$ 为本征形式。则存在有限扩张 $K'|K$、位扩张 $\p'|\p$ 以及本征形式
\[
\tilde g\in M_k\bigl(\Gamma_1(N);\OO_{K',(\p')}\bigr),
\]
使得 $\tilde g$ 模 $\p'$ 约化为 $g$。

\noindent\textbf{证明}。
\begin{itemize}
\item 考虑环同态
\[
\mathbb{T}_{\OO_K}\longrightarrow\mathbb{T}_{\F_{\p}}\xrightarrow{\ \psi_g\ }\F_{\p}.
\]
\item 这对应一个极大理想 $\m$。由于 $\Spec(\mathbb{T}_{\OO_K})\to\Spec(\OO_K)$ 局部自由，可以找到素理想 $\mathfrak{q}\subseteq\m$，使得 $\mathfrak{q}\cap\OO_K=(0)$。
% ===== PDF p.083 =====
\item 这定义了 $\mathbb{T}_K$ 的一个非零素理想，对应 Hecke 特征系统 $\mathbb{T}_K\to K'$。
\item 由引理 10.9，这对应一个本征形式。
\item 用引理 10.9 的唯一性可知它约化为 $g$。\hfill$\square$
\end{itemize}
% ===== PDF p.084 =====
\subsection*{定理 10.8 的证明}

\begin{itemize}
\item 设 $f$ 是权 $\ge 2$ 的模 $p$ 本征形式，定义在 $\F_p$ 的扩域 $F_f$ 上。
\item 把它提升为某数域 $K$ 上的本征形式 $\tilde f$；设 $\p$ 为 $K$ 中位于 $p$ 之上的位。
\item 与之相伴，有 Galois 表示
\[
\rho_{\tilde f}:G_\Q\longrightarrow\GL_2(K_{\p}).
\]
\item 由于像紧，总可以找到同构的表示
\[
\rho_{\tilde f}:G_\Q\longrightarrow\GL_2(\OO_{K,\p}).
\]
\item 把它模 $\p$ 约化，得到表示
\[
\rho_{\tilde f}:G_\Q\longrightarrow\GL_2(\F_{\p}).
\]
\item 它未必已经半单。于是把 $\rho_f$ 定义为半单化
\[
\rho_f:=\rho_{\tilde f}^{\mathrm{ss}}
\]
（即 Jordan--Hölder 因子的直和）。它与 $\rho_{\tilde f}$ 有相同的迹与行列式。
\item 由于 $\ker\rho_f\subseteq\ker\rho_{\tilde f}$，它在 $pN$ 外非分歧。
\item 还需说明：$\rho_f$ 已在 $F_f$ 上定义。为此只需证明：对一切 $\sigma\in\Gal(\overline{\F}_p|F_f)$，有 $\rho_f=\sigma\circ\rho_f$。
\item 当然，对一切 $\ell\nmid pN$，$\Frob_\ell$ 的特征多项式
\[
X^2-a_\ell X+\varepsilon(\ell)\ell^{k-1}
\]
都被 $\sigma$ 保持，因为按定义 $a_\ell,\varepsilon(\ell)\in F_f$。现在用：\znoteplain{原文此处写作 $X^2-a_\ell X+\varepsilon(\ell)$，漏掉了权 $k$ 的因子 $\ell^{k-1}$；由 $\det\rho_f=\varepsilon\chi_{\mathrm{cycl}}^{k-1}$ 可知应为 $X^2-a_\ell X+\varepsilon(\ell)\ell^{k-1}$。译文补全。}
\end{itemize}

\noindent\textbf{定理 10.11}（Brauer--Nesbitt）\phantomsection\label{thm:10.11}\quad 设 $F$ 是完美域，$G$ 是有限群，$V,W$ 是 $G$ 在 $F$ 上的两个半单有限维表示。则 $V\cong W$ 当且仅当对一切 $g\in G$ 都有
\[
\operatorname{charpoly}(g\mid V)=\operatorname{charpoly}(g\mid W).
\]

\begin{itemize}
\item 当系数域的特征为正时，这取代了第 6 节中关于迹的类似陈述。
\item 至此完成了「$f\rightsquigarrow\rho_f$」的构造。\hfill$\square$
\end{itemize}

\subsection*{Serre 猜想\idx{Serre 猜想}}

\begin{itemize}
\item 在继续讨论权 $1$ 形式之前，先作一段简短的插叙。
\item 1973 年，Serre 提出了一个猜想，恰好刻画了哪些模 $p$ Galois 表示以这种方式出现。它被称为
\end{itemize}

\noindent\textbf{猜想 10.12}（Serre 模性猜想，或称 Serre 猜想）\phantomsection\label{conj:10.12}\quad 设 $\rho:G_\Q\to\GL_2(\F_p)$ 是连续、奇、不可约的 Galois 表示。则存在模 $p$ 本征形式 $g$，使得
\[
\rho=\rho_g.
\]

1987 年，Serre 发表了一篇文章，细化了这个猜想：

\begin{itemize}
\item 他预测了 $g$ 的最优级（用 $\rho$ 的 Artin 导子\idx{导子}表出）。
\item 他预测了最优权（用 $\rho$ 在 $p$ 处的分歧给出配方）。
\item 他证明了：他的猜想将蕴含 Fermat 大定理。
\item Serre 猜想现已成为定理，由 Khare--Wintenberger（2008）证明。
\end{itemize}
% ===== PDF p.085 =====
\subsection*{把 $\rho_f$ 提升到 $\C$}

\begin{itemize}
\item 回到先前的设定：$f\in M_1(\Gamma_0(N),\varepsilon)$ 是本征形式。
\item $K$ 是由特征值与 $\varepsilon$ 生成的数域。
\item 令 $L$ 为在 $K$ 中完全分解的 $\Q$ 的素数之集。由 Chebotarev 密度定理，这是无限集。
\item 设 $\ell\in L$，$\lambda$ 为 $K$ 中位于 $\ell$ 之上的位，$\OO_{K,(\lambda)}\subseteq K$ 为 $\lambda$-整元素。
\item 现在把上一节应用于 $p$ 由变动中的 $\ell\in L$ 扮演的情形。
\item 对每个这样的 $\lambda$，考虑 $f$ 模 $\lambda$ 的约化 $\bar f$；这是模 $\ell$ 的本征形式。
\item 定理 10.8 给出半单 Galois 表示
\[
\rho_{f,\ell}:G_\Q\longrightarrow\GL_2(\F_\ell)
\]
（因 $\lambda$ 完全分裂，系数域就是 $\F_\ell$）；对任一 $p\nmid\ell N$，$\Frob_p$ 的特征多项式为
\[
\operatorname{charpoly}(\Frob_p\mid\rho_{f,\ell})=X^2-a_pX+\varepsilon(p).
\]
\item 它的零点是 $n$ 次单位根；由下一引理，可以假设它们全在 $\OO_K$ 中，从而在 $\F_\ell$ 中。于是该特征多项式 $=(X-a)(X-b)$，其中 $a,b\in\F_\ell^\times$。\znoteplain{这里用了两件事：其一，$K$ 已扩大为包含所有 $n\le A$ 次单位根，故 $\F_\ell^\times$ 中的 $n$ 次单位根在 $K$ 中的提升（Teichmüller 提升）落在 $K$ 内，且是 $\lambda$-整的，即在 $\OO_{K,(\lambda)}$ 中；其二，$\lambda$ 完全分裂，故剩余域就是 $\F_\ell$，这些单位根模 $\lambda$ 后落在 $\F_\ell$ 中。}
\end{itemize}

\noindent\textbf{定义 10.13}\phantomsection\label{def:10.13}\quad 令 $G_\ell:=\operatorname{im}(\rho_{f,\ell})\subseteq\GL_2(\F_\ell)$。这是有限群。

\noindent\textbf{引理 10.14}\phantomsection\label{lem:10.14}\quad 若 $f$ 的权为 $1$，则存在 $A\in\N$，使得对一切 $\ell\in L$ 都有 $|G_\ell|\le A$。

\noindent\textbf{证明}\quad 在用到关于权 $1$ 模形式系数的一个关键估计之后，这就变成纯群论问题：对 $\GL_2(\F_\ell)$ 的子群分类。我们不讨论它；若想看，见 \texttt{[DS74, Lemma 8.4]}。\hfill$\square$

\begin{itemize}
\item 扩大 $K$，使其包含所有 $n\le A$ 的 $n$ 次单位根。这会缩小 $L$。
\item 再把 $\le A$ 的素数从 $L$ 中排除；此时 $L$ 仍无限。
\end{itemize}

\noindent\textbf{引理 10.15}\phantomsection\label{lem:10.15}\quad 设 $G$ 是阶与 $\ell$ 互素的有限群。则任一表示 $\rho:G\to\GL_2(\F_\ell)$ 都可提升为表示
\[
\tilde\rho:G\longrightarrow\GL_2(\OO_{K,(\lambda)}).
\]

\noindent\textbf{证明}。\textbf{第一步：提升到 $\GL_2(\Z_\ell)$}。

\begin{itemize}
\item 只需证明：任一 $G\to\GL_2(\Z/\ell^n\Z)$ 都可提升为 $G\to\GL_2(\Z/\ell^{n+1}\Z)$。
\item 为此考虑短正合列
\[
1\longrightarrow 1+\ell^nM_2(\F_\ell)\longrightarrow\GL_2(\Z/\ell^{n+1}\Z)\longrightarrow\GL_2(\Z/\ell^n\Z)\longrightarrow 1.
\]
\item 第一个群通过 $1+\ell^nx\mapsto x$ 同构于 $M_2(\F_\ell)$。
\item 现在用非交换群上同调：考虑左作用，$G$ 通过
\[
(g\cdot\phi)(x):=\phi(g)^{-1}\phi(gx)
\]
作用在 $\mathrm{Map}\bigl(G,\GL_2(\Z/\ell^n\Z)\bigr)$ 上。\znoteplain{原文此处在提取文本中作 $g\cdot\phi:=\phi(g)\phi(g)^{-1}$，明显残缺。按非交换群上同调的标准扭作用应为 $(g\cdot\phi)(x)=\phi(g)^{-1}\phi(gx)$：其不动点恰为同态 $\phi:G\to\GL_2(\Z/\ell^n\Z)$，故有下一个等式。译文据此重建。}
\item 于是
\[
\Hom\bigl(G,\GL_2(\Z/\ell^n\Z)\bigr)=\mathrm{Map}\bigl(G,\GL_2(\Z/\ell^n\Z)\bigr)^G.
\]
\item 对上述短正合列施加 $\mathrm{Map}$，再由 $G$ 得到非交换群上同调的长正合列
\[
\cdots\to\Hom(G,\GL_2(\Z/\ell^{n+1}))\to\Hom(G,\GL_2(\Z/\ell^n))\to H^1\bigl(G,\mathrm{Map}(G,M_2(\F_\ell))\bigr).
\]
\item 群上同调的事实：若 $|G|$ 与 $|A|$ 有限且互素，则 $H^1(G,A)=0$。
\end{itemize}
% ===== PDF p.086 =====
\noindent\textbf{第二步：已在数域上定义。}

\begin{itemize}
\item 由于 $G$ 有限，$\rho$ 半单。
\item 迹必为阶 $\le|G|\le A$ 的单位根之和，从而落在 $K$ 中。
\item 由 Brauer--Nesbitt 可知 $\rho$ 已在 $K$ 上定义。
\item 最后，$K\cap\overline{\Z}_\ell=\OO_{K,(\lambda)}$。\hfill$\square$
\end{itemize}

\subsection*{回到定理 10.2 的证明（权 1 形式的 Galois 表示）}

\begin{itemize}
\item 把它应用于 $\rho_{f,\ell}:G_\ell\hookrightarrow\GL_2(\F_\ell)$，得到 $G_\ell\to\GL_2(\OO_{K,(\lambda)})$。\znoteplain{原文这一段的提取文本把下标时而写作 $g$、时而写作 $f$（例如 $\rho_{g,\ell}$、$\tilde\rho_{g,\ell}$ 与 $\rho_{f,\ell}$ 并存）。按上下文，本讲自始至终讨论的是同一个权 $1$ 本征形式 $f$，故译文统一作下标 $f$。}
\item 再与 $G_\Q\to G_\ell$ 复合，得到 Galois 表示
\[
\tilde\rho_{f,\ell}:G_\Q\longrightarrow\GL_2(\OO_{K,(\lambda)}).
\]
\end{itemize}

\noindent\textbf{命题 10.16}\phantomsection\label{prop:10.16}\quad 提升 $\tilde\rho_{f,\ell}:G_\Q\to\GL_2(\OO_{K,(\lambda)})$ 满足：对一切 $p\nmid\ell N$，
\[
\operatorname{charpoly}(\Frob_p\mid\tilde\rho_{f,\ell})=X^2-a_pX+\varepsilon(p).
\]

关键在于变动 $\ell$：回忆我们对所有 $\ell\in L$ 都得到 $\tilde\rho_{f,\ell}$。

\noindent\textbf{定义 10.17}\phantomsection\label{def:10.17}\quad 令 $Y$ 为形如 $(X-\alpha)(X-\beta)$ 的多项式之集，其中 $\alpha,\beta\in\OO_{K,(\lambda)}$ 是阶 $\le A$ 的单位根。这是有限集。

\begin{itemize}
\item \textbf{注}：给定 $n\le A$、$\ell\in L$ 以及 $n$ 次单位根 $x\in\F_\ell^\times$，存在唯一的提升为 $K$ 中的单位根（存在性：$n\le A$；唯一性：$\ell\ge A$）。
\item \textbf{推论}：若 $F,G\in Y$ 且 $F\equiv G\pmod\lambda$，则 $F=G$。
\item \textbf{想法}：把它应用于 $F=X^2-a_pX+\varepsilon(p)$ 与 $G=\operatorname{charpoly}(\Frob_p\mid\tilde\rho_{f,\ell})$。
\end{itemize}

\noindent\textbf{引理 10.18}\phantomsection\label{lem:10.18}\quad $X^2-a_pX+\varepsilon(p)\in Y$。

\noindent\textbf{证明}\quad 对每个位于 $L$ 中素数之上的 $\lambda$，由于 $\rho_{f,\ell}$ 的像有限，存在 $F\in Y$ 使得
\[
\operatorname{charpoly}(\Frob_p\mid\rho_{f,\ell})\equiv F\pmod\lambda.
\]
于是也有
\[
X^2-a_pX+\varepsilon(p)\equiv F\pmod\lambda.
\]
因 $Y$ 有限，存在某个 $F$ 使上式对无穷多个 $\lambda$ 成立。故 $X^2-a_pX+\varepsilon(p)=F\in Y$。\hfill$\square$

\noindent\textbf{命题的证明}。
\begin{itemize}
\item 由于 $G_\ell$ 有限且阶 $\le A$，有 $\operatorname{charpoly}(\Frob_p\mid\tilde\rho_{f,\ell})\in Y$。
\item 又
\[
\operatorname{charpoly}(\Frob_p\mid\tilde\rho_{f,\ell})\equiv\operatorname{charpoly}(\Frob_p\mid\rho_{f,\ell})\equiv X^2-a_pX+\varepsilon(p)\pmod\lambda.
\]
\item 两边都在 $Y$ 中，故 $\operatorname{charpoly}(\Frob_p\mid\tilde\rho_{f,\ell})=X^2-a_pX+\varepsilon(p)$。\hfill$\square$
\end{itemize}

尚需证明：

\begin{itemize}
\item $\rho$ 为奇：这仅因为 $\varepsilon$ 为奇且 $c$ 的阶为 $2$。
\item $\rho$ 不可约：用 Rankin 的一个复分析估计。\hfill$\square$
\end{itemize}

至此完成了「本征形式 $\rightsquigarrow$ Galois 表示」的构造！
% ===== PDF p.087 =====
\subsection*{权 1 构造的总结}

\begin{center}
\begin{tikzcd}[row sep=3.4em, column sep=2.6em]
\text{权 }1\text{ 本征形式}
  \arrow[r, "{\text{同余}}"]
  \arrow[d, "{\text{模 }\ell\text{ 约化}}"']
& \text{权 }\ge 2\text{ 本征形式}
  \arrow[r, "{\text{Deligne 构造}}"]
  \arrow[d, "{\text{模 }\ell\text{ 约化}}"']
& \ell\text{-进 Galois 表示}
  \arrow[d, "{\text{模 }\ell\text{ 约化}}"]
&
\\
\text{模 }\ell\text{ 的权 }1\text{ 本征形式}
  \arrow[r, "{\cdot\,\mathrm{Ha}}"]
& \text{模 }\ell\text{ 的权 }\ge 2\text{ 本征形式}
  \arrow[u, "{\text{DS 提升}}"]
  \arrow[r, "{\text{模 }\ell\text{ 构造（定理 10.8）}}"]
& \text{模 }\ell\text{ 的 Galois 表示}
  \arrow[r, "{\text{提升}}"]
& \text{复 Galois 表示}
\end{tikzcd}
\end{center}

\znoteplain{原第 87 页是一张示意图（方框加箭头）。此处按原图的七个方框重排为两行四列的交换图：「$\cdot\,\mathrm{Ha}$」表示乘以 Hasse 不变量，「DS 提升」即引理 10.10，「模 $\ell$ 构造」用定理 10.8，最后的「提升」即本讲末尾把 $\rho_{f,\ell}$ 提升到 $\C$ 的步骤。原图为扫描排版，箭头走向按图中文字标注复原，宜与原书对照。}
